\global\nosectionbreakfalse
\phantomsection\addcontentsline{toc}{section}{Общие указания к выполнению лабораторных работ}
\begin{center}\textbf{ОБЩИЕ УКАЗАНИЯ К ВЫПОЛНЕНИЮ ЛАБОРАТОРНЫХ РАБОТ}\end{center}
\vspace{4pt}

\labpart{Порядок выполнения и защиты}
Перед выполнением работы изучают указанные в ней разделы теории и получают номер варианта. На защите
студент демонстрирует работу программы, изменяя параметры по указанию преподавателя, и отвечает на
контрольные вопросы.

\labpart{Программная среда}
Работы выполняются на Python~3 в Jupyter: NumPy, Matplotlib или Plotly, SciPy (подраздел~\ref{sec:python},
приложение~\ref{app:python}). Все алгоритмы реализуются \emph{самостоятельно} на NumPy: это требует
понимания каждой формулы и позволяет исследовать внутренние величины алгоритма. Готовые библиотеки
(scikit-learn, PyTorch, MiniSom, DEAP, PyGAD) служат для сравнения, расхождение с ними нужно объяснить.
MATLAB и JavaScript допускаются по желанию студента.

\labpart{Требования к программе}
\begin{itemize}
\item параметры алгоритма задаются в интерфейсе (\code{ipywidgets}) или в одном блоке в начале программы;
\item процесс обучения или эволюции отображается динамически, с заданным периодом перерисовки;
\item генератор случайных чисел создаётся с фиксированным seed (\code{np.random.default\_rng(seed)});
\item для стохастических алгоритмов предусматривается серия из $M$ независимых запусков со сводной
статистикой: доля успехов, медиана и худший результат, сравнение с эталоном
(подраздел~\ref{sec:ga-method}).
\end{itemize}

\labpart{Содержание отчёта}
Отчёт по СТП БГУИР содержит: титульный лист, цель, постановку задачи и вариант, краткую теорию,
описание реализации, результаты с графиками, таблицами и их анализом, выводы, опирающиеся на числа,
ответы на контрольные вопросы и листинг программы в приложении.
